% -----------------------------------------------------------
% -----------------------------------------------------------
\part{Blessings, Lessons I've Learned, and Spiritual Experiences}
% -----------------------------------------------------------
% -----------------------------------------------------------
	\label{LESSONS}
	
% -----------------------------------------------------------
\section{Blessings} % *BLESSINGS
% -----------------------------------------------------------
	\label{BLESSINGS}
	
	% ----------------------
	\subsection{Patriarchal Blessing (3 July 2005)}
	% ----------------------
		\label{BLESSINGS-PATRIARCHAL}
		
		% -----
		\subsubsection{Introduction}
		% -----
		
			I received this blessing on July 3, 2005. It was given by George Richard Runyan, who at the time was the patriarch of the Vernal Utah Ashley Stake. I received it at his home, with both of my parents present. 

			The sentence numbers are not part of the original blessing; I have added these to simplify references I make to the blessing.
		
		% -----
		\subsubsection{Blessing}
		% -----
		
			1. Bryce Scott Gessell by the authority of the holy Melchizedek priesthood and as a patriarch I lay my hands upon your head and give you your patriarchal blessing.
			2. Bryce your Father in Heaven loves you and is very pleased with you and your desires to serve him.
			3. You are of the tribe of Ephraim and through that tribe you will receive your blessings.
			4. The blessings of Abraham Isaac and Jacob are very important to all who have accepted the gospel.
			5. As a young man you will have the responsibility to go into the world and teach the gospel of Jesus Christ, to bring souls unto Christ and help prepare the world for the coming of the Savior.
			6. As a missionary you will have the Spirit to guide and direct you.
			7. I bless you that you will be a good missionary.
			8. Obey the mission rules and you will have success in your work and in your assignments.

			\textbf{[paragraph break]}

			9. Your example and testimony will make it possible for you to be a leader in the Church.
			10. I bless you that you will have the keys of the priesthood at one time or another to administer in the kingdom and to lead the work.
			11. In your leadership callings I bless you that you will have patience.
			12. The ability to learn from others.
			13. Take counsel and make decisions which are appropriate for the circumstances at the time.

			\textbf{[paragraph break]}

			14. I would admonish you to always keep your mind focused on those things which are eternal and not let worldly things become important in your life.
			15. As a servant of your Father in Heaven you will be able to affect the lives of many people.
			16. They will love and appreciate you for your service.
			17. Remember you are a child of your Father in Heaven and a servant of Jesus Christ.
			18. Strive to be like Him and follow His example.

			\textbf{[paragraph break]}

			19. As you determine what your life occupation will be I bless you that you will be directed by the Spirit.
			20. You will find something that will provide for you and your family.
			21. Keep in mind that your availability to your Father in Heaven is very important in giving service.

			\textbf{[paragraph break]}

			22. I bless you that you will be able to find a wonderful and sweet woman to take to the Temple and be sealed for time and all eternity.
			23. One you can live with congenially and raise a family and have the Spirit in your home.
			24. I bless you that you will have the ability to lead your family as a righteous father and be the example that they need to help them become what they should be.

			\textbf{[paragraph break]}

			25. You will have the opportunity to be a leader in the community.
			26. Through this experience you will have the opportunity to bear witness of the gospel truths and teach the gospel.
			27. There will be many who will be affected by your example and service.

			\textbf{[paragraph break]}

			28. I bless you with the gift of healing.
			29. Through the Melchizedek priesthood and the Spirit you can administer and pronounce blessings that would be appropriate, which will benefit your family and those with whom you work.

			\textbf{[paragraph break]}

			30. I bless you with the ability to find those who are in need.
			31. You will be able to determine what they lack and how you can serve them.
			32. You will be able to help people who are less active and those who do not know the gospel.

			\textbf{[paragraph break]}

			33. I bless you that you will be successful in your occupation.
			34. Remember to seek first the kingdom of God and the things which you need will be supplied for you.

			\textbf{[paragraph break]}

			35. Make your home and your life a place where the Spirit can dwell.
			36. Do not allow into your home anything that would drive away the Spirit.
			37. There is so much evil in the world.
			38. You need to protect yourself and your family from Satan's influence.
			39. As a husband, a father and a priesthood holder you can bring into your home that which would help each one love the gospel and return to live with their Father in Heaven.

			\textbf{[paragraph break]}

			40. You will have difficulties and trials in your life.
			41. Your strong testimony and love for your Father in Heaven and your Savior Jesus Christ will carry you through any difficulties that you might have.

			\textbf{[paragraph break]}

			42. I bless you with a quickness of mind that you be able to learn quickly and retain those things which you need to know.
			43. You will be able to be efficient in your service to your Father in Heaven and in your occupation.

			\textbf{[paragraph break]}

			44. You will have the gift of understanding to be able to evaluate circumstances and make good decisions.
			45. Those who are around you will benefit from this blessing.

			\textbf{[paragraph break]}

			46. These blessings are all contingent upon your faithfulness and through the will of the Lord.
			47. In the name of Jesus Christ, amen.
			
	% ----------------------
	\subsection{Blessing from President Openshaw (16 June 2007)}
	% ----------------------
		\label{BLESSINGS-OPENSHAW}
		
		% -----
		\subsubsection{Introduction}
		% -----
		
			I received this blessing on June 16, 2007. President Openshaw finished his service a few weeks later and went home. He gave a personal blessing to every departing missionary during that missionary's final interview. Since I knew he would be gone by the time I had my final interview, I asked him to give me an interview early, which he did. This blessing was the final part of the interview. 

			Before he gave the blessing, I prayed so that I would be able to remember as much as possible from what he said. As soon as he was done, I thanked him and went to a large upstairs bathroom, where I prayed again and then wrote as much as I still had in my mind. I no longer know if I missed anything; my impression at the time was that I had missed one or two statements but had recorded everything else. I have no way of checking the accuracy of what I have, but I did not write anything that I didn't remember distinctly. 
						
			I added the sentence numbers later.
		
		% -----
		\subsubsection{Blessing}
		% -----
		
			1. The Lord recognizes the sacrifice that you have made not only in your time and in your life but in opening up your heart.

			2. I bless you so that this process will continue.

			3. You will have many callings---important callings---in the church which will be of grave importance in the lives of others and the building of the kingdom.

			4. I bless you with the gift of discernment, to see the needs of others and how to resolve them, and not only that, but to see beyond what they want to hear in order to see what the Lord wants them to know.

			5. I bless you with the courage to tell it to them.

			6. There will be people who will be offended by your honesty but over time a reputation of your integrity and loyalty to the Lord will grow and there will be those who will come from other stakes and other parts of the world to seek your guidance.

			7. There will be moments in your life in which the temptation to leave the church behind along with the testimony you have gained and the covenants you have made will be strong, but I bless you that you will never, never fall away from what you know and the testimony you have gained.

			8. I bless you that you will do the four things that Alma taught Helaman his son---pray, read the scriptures, live the life of a child of God, and bear testimony of what you know.

			9. These things will be very important.

			10. I bless you to never forget the experiences that you have had in your mission, and to remember the profound feelings that have welled up in you in lessons with converts.\footnote{``Welled up in you'' is not the exact original phrase, but I have forgotten what it was.}

			11. I bless you to not only remember the experiences with converts, but the experiences with your companions (you will have many trials in your life, but the beauty of their lives will be a source of comfort for you).

			12. Communicate with them.

			13. I bless you to seek a university in which you will be able to testify of Jesus Christ and in which you can bear witness of what you know.

			14. I bless you to have a wife that loves you.

			15. I bless you that you will have the humility you need to accept the counsels of your wife and to apply them in your life.

			16. I bless you that through daily routines and habits, you will have the power to overcome the temptations which you face in your life.

			17. Through habits you will be able to develop yourself.

			18. I bless your children to follow your example.

			19. They will love the scriptures because you love the scriptures.

			20. They will talk of doctrine because you talk of doctrine.

			21. They will want to serve others because they will see you and your wife doing it, and they will be good missionaries.

			22. I bless you so that the Lord can be as close to you as you would like him to be.

			23. I bless you to not only feel the love of the Lord, but to feel His presence in your heart, and to know He is there.

			24. I bless you to feel His presence very closely in the time between your arrival at home and the beginning of your career.

			25. I bless you that your heart will always follow the promptings of the Holy Spirit.

			26. I bless you that you will be able to work longer and harder than most people.

			27. I bless your mind to be open to new thoughts and new ways of thinking, that you will be able to study and understand, and that you will be able to analyze material and return it with creativity and skill.

			28. I bless your body from your feet to your head, so your muscles and bones and organs will work as they were designed to work, and they will respond to the priesthood.

			29. I bless you with health, to avoid sickness, and to be healthy.

			30. I bless you to be able to remember the things you study.

			31. (President Openshaw, at a gathering of returned missionaries, October 5, 2007: (``The only way you aren't going to receive these blessings if is you do not accept them.'')

	% ----------------------
	\subsection{Blessing from Scott Gessell (5 August 2018)}
	% ----------------------
		\label{BLESSINGS-2018}
		
		% -----
		\subsubsection{Introduction}
		% -----
		
			I don't remember much about the circumstances of this blessing. I was in Utah at the time because we were living in Provo while I taught at BYU for a summer term. Being at BYU had been difficult, but we were having a good time overall, I think. 
			
			I think what happened is that I was having a very difficult time gearing up for the job market in the fall, and was going through the extreme uncertainty and anxiety that went along with that. I didn't know where I stood with BYU---and, as I found out later, my standing with them wasn't good---and I was worried about other possible jobs as well. 
			
			From the few journal entries I had at that time, it's clear I was feeling very worried:
			
				\begin{compactdesc}
					\item[30 June 2018] I went to the library briefly to look over some materials for class on Monday, since I haven't taught this stuff before (Newton), and Travis Anderson said he was going to try to observe, so I feel like I need to be more prepared than normal. It was not awkward meeting Travis. I hope he is able to come and I've been praying a lot that things will go well.
					\item[16 July 2018] Today was a good day for me at work. I talked to Felipe and got some clarity about what I've been working on, and I got back on track. I feel better, but it's been hard. This is not my ideal situation. I have liked BYU and have felt more comfortable than I thought I would, but this still is hard for me. Less routine, not at our own house, all of that. Not my strength. 
					\item[7 August 2018] Man I have been really bad at writing. It's a shame too because we have had a lot of fun. So much seeing family, so much with the kids, Vernal, it's just been great. Extremely difficult for me for lots of reasons, but great.
				\end{compactdesc}
				
			I asked my dad for this blessing just a few days before writing that entry for 7 August, and so I was really messed up with everything that was happening. Oh, and Alicia was pregnant with Felicity.
			
			That's all the background I know, at least for now.
			
		% -----
		\subsubsection{Blessing}
		% -----

			Bryce Scott Gessell, in the name of Jesus Christ, and by the authority of the Melchizedek priesthood, I lay my hands upon your head, to give you this blessing of comfort and of counsel, and of direction, and indeed knowledge in your life.

			Bryce, you have been faithful and diligent and loyal to the things that are so important to you and to your family; to the gospel, to your future career, to your schooling. I bless you that you will continue to be diligent and loyal and active in the things that are so important to you, in these things from God, that has given you great blessings in your life. 

			Remember the scripture, Doctrine and Covenants 82:3, that says, where much is given, much is required. You have been given much in your life, not only because of your membership in the gospel of Jesus Christ, but because of your loyalty, and the things that you have been doing and the choices you have made to this point. I bless you that you will continue to be wise, that you will continue to make decisions that are beneficial for yourself and especially your family.

			At this point in your life, Bryce, you are coming upon a crossroads, and a culmination of many years of work; and decisions will need to be made in the future that will have a great impact upon you and your family. I bless you that you will stay extremely close to the Spirit, and that you will invite personal revelation into your life. Seek after this personal revelation diligently through your scripture study, through personal prayers, and through temple worship. 

			Remember that you, Bryce, are the patriarch of your family; and because of this, you will feel anxious about many of the decisions that you have to make. But remember that Heavenly Father does not want you to make these decisions alone. You need to include him, in your personal prayers and in your pleadings, and you also need to include Alicia, as she has been the foundation for many of your blessings. Counsel together often; see that the needs of your family are met. There are so many different variables that come into play, and that will come into play in the coming years---do not take on this responsibility of decision-making alone, but include others in your life. 

			I bless you that you will seek out the support beyond our Heavenly Father and beyond Alicia, that you will seek out the support of those who will make decisions on your behalf. Use your personality, Bryce, use your ability to communicate with them. Be up front and open and honest with them; talk to them often, and remember that opportunities that Heavenly Father provides are for your benefit; that you have been given the right to choose which opportunities are best for you and your family. 

			I further bless you that you will remain calm, you will remain collected, and you will be focused on those things which are beneficial first and foremost to your family, and not just in the short term, but in the long term. 

			I bless you that you will have the desire and the energy you need to go above and beyond expectations of others, and even beyond your own expectations, that you will find success in all your endeavours. Remember that, as you move forward in life, there will be times when challenges and difficulties arise, that will seem to put roadblocks and detours in front of you. Know that these are only temporary impediments to your future, and you have been given the gift of knowledge and education, in order to overcome anything that Satan puts in your way. 

			I bless you that you will continue to remain calm with the baby, the upcoming addition to your family; that you will be able to adapt well to a third child, and that you will give that child the love that is needed in order to progress here on this earthly life. 

			Now Bryce as you move forward, remember that you have a particular personality and teaching style, and it is of utmost importance that you be open with your teaching style, and that you find a place to teach that will be well-suited to your style. 

			And finally Bryce, I bless you that you will be able to look forward to the things that are in store for you, knowing that so many doors will open for you. I bless you that you will be excited about the future and the possibilities, and the things that haven't happened yet. 

			Bryce you are a wonderful father, you are a wonderful husband; may you continue on on your journey in this mortality, knowing that Heavenly Father loves you, knowing that your parents love you, and that Alicia, Addy, and Brigham love you so much. Heavenly Father is indeed proud of you as one of his sons, and I leave this with you, by the authority of the Melchizedek priesthood and in the name of our lord and our savior Jesus Christ, amen.

% -----------------------------------------------------------
\section{Spiritual Experiences}
% -----------------------------------------------------------
	\label{SE}
	
	% ----------------------
	\subsection{Teaching an impromptu lesson in the Rockbridge 1st ward in 2024}
	% ----------------------
		\label{SE-LESSON-CTR}
		
		(See my journal entry for 17 November 2024.)
		
		The Spirit led me to be prepared for this lesson without realizing it. A few weeks ago I was reading in Ether and noticed something that I wanted to talk about with Climbing the Rainbows, so I adjusted my schedule so that I could write about these chapters for this week's essay, which came out this Friday. And then I got asked to teach the lesson today. I did not experience any of these thoughts as promptings from the Spirit but that is what happened. I do wish that I were teaching Sunday school, since I enjoy that calling the most of any calling I've had.

	% ----------------------
	\subsection{Giving the ``Why I Believe'' devotional at the institute in 2023}
	% ----------------------
		\label{SE-WHY}
		
		Written 10/6/23
		
		In fall 2023 I was asked by the institute council to give a devotional talk on the topic ``Why I Believe.'' I had a bit under two weeks to prepare and I initially wrote down some ideas in the first few days after they asked me but I didn't pick it up again until the Sunday before I was supposed to speak (on a Friday).
		
		I struggled to decide what to speak about. I wrote out many pages of notes on different topics, but as I went nothing really felt right. I wrote out some more notes on a different occasion and it still felt off. I was feeling lost, so I decided to fast on the Wednesday before I was to speak. I fasted and during the morning it occurred to me that I should just ask the students what they thought about it and what they wanted to hear. I did that in the minds/brains class and I ran into Quin and Kamryn Meyer on campus and asked for their thoughts too. They all gave me some very useful thoughts. I wrote them all down and later that same day I started writing down some new ideas that I felt were being given by the Spirit, and I continued to write during faculty meeting on Wednesday night.
		
		On Thursday, the day before I spoke, I worked for a while more and got most of the new ideas in order. I had a late meeting that night and I went to it and when I got home I gave a practice talk to Alicia that had most of the material I felt like I had received after speaking to the students and fasting.
		
		Alicia didn't really like it. She said it had no structure and didn't really answer the question of why I believe. We talked it over some and I agreed with some of the things she said and started to write out some different things in a different order. But I was really frustrated because I was really tired of working on it and tired of worrying about it, and also just tired because I wanted to go to bed. 
		
		The next morning, Friday, I woke up and grabbed my computer and typed a few additional ideas but I mainly prayed for the Spirit. I said, I still don't really feel good about any of this, so please give me the Spirit in the moment so I can select from among all the ideas I've prepared.
		
		The devotional came around and I got up to speak. I forgot to record it, which was a shame. But I started talking and only glanced at my notes a couple of times. What I said ended up being a mix of everything I had written at various times, and there was also material I hadn't written.
		
		People were receptive to what I said and it resonated with some of them. In conversations after the talk, I said a number of other things that I hadn't said in the talk, but which I had prepared.
		
		I had class soon after the devotional was over. I hadn't prepared at all because I was so focused on the talk, but in that class we talked about my talk and about many other related issues. I ended up saying many more things that I had written down during my preparation, and some of those things were fairly direct answers to questions that students asked me during class.
		
		Preparing for this devotional was a huge, frustrating struggle. It felt like I just couldn't get anything right and everything felt off. I went into the speech not feeling good about almost any of it. But I realize now that the Lord was preparing me to speak not only in that devotional, but to people after the devotional and then again to my students in class. I needed to be ready to share \textit{all} that material on different occasions, and the things I shared came from every prep session and version of the talk that I made. If I had gotten things right the first time, I wouldn't have been able to prepare all the other material that I ended up using outside the devotional. I needed all of that to talk to everyone throughout the day.
		
		In addition, several students emailed me after the devotional wanting to meet to talk over the things I had spoken about. 
		
		I felt that the Lord had me struggle through prep for a larger purpose that I was not aware of until after it had happened, and I realized that I hadn't felt right about the prep because I needed to get a bunch of other things ready.
		
		Another important lesson here is how the Lord used uncomfortable and unsettling feelings in order to get me to act in a particular way. It was the opposite of peace---it was not pleasant, it was anxiety-inducing, and it literally kept me up and night. But there was a job to do and the job got done, and unsettling feelings were the way to accomplish it.
		
		It was also a lesson in preparing as well as I can and then turning it over to the Lord by asking the Spirit. He provided it and guided me during the talk.

% -----------------------------------------------------------
\section{Lessons about Myself}
% -----------------------------------------------------------
	\label{LESSONS-ME}
	
	% ----------------------
	\subsection{My professional life}
	% ----------------------
	
		% -----
		\subsubsection{Meetings}
		% -----
		
			\begin{itemize}
			
				\item \textbf{Run every meeting, no matter what, as though each person there is waiting for it to end so that they can go take care of young children.} Don't drag it out; when it's over, end it; begin on time; end on time; get through it.
				
				\item \textbf{JUST BECAUSE YOU ARE A LEADER DOES NOT MEAN PEOPLE WANT OR NEED TO HEAR FROM YOU!} You do not need to finish every meeting by talking. You are not that interesting. You are not that insightful. Stop talking. Stop talking! STOP TALKING!!!
			
			\end{itemize}
	
		% -----
		\subsubsection{Research}
		% -----
	
			\begin{itemize}
			
				\item \textbf{I am never going to stop wanting to research and write in contemporary philosophy. There will be times when I'm less interested, but I will always want to come back to it, so always leave the door open and plan for my desire to return.}
			
			\end{itemize}
			
		% -----
		\subsubsection{Job searching and job market}
		% -----
		
			\begin{compactitem}
			
				\item \textbf{The academic job market is \textit{much} more work than you realize or remember! Unless you are actively job-searching or have just completed an attempt at the job market, you don't remember what it's like! Do not decide flippantly to job-search in academia!}
					\begin{compactitem}
						\item The work of preparing job documents, deciding where to apply, writing cover letters and customizing job materials, and so on, is \textit{much more time-consuming than you remember}.
						\item This work is also \textit{much more draining than you remember}. If you go on the market all the way, you will \textit{not} have energy to work on other writing or research projects.
						\item You will have to set aside all your other projects until you are completely done filling out \textit{all} your materials, and even then, you may have to prepare for interviews, which will \textit{also} take a long time.
						\item Do \textit{not} think that you can stack applying for jobs onto everything else! You cannot! You will not get anything else done!
					\end{compactitem}
			
			\end{compactitem}

% -----------------------------------------------------------
\section{Lessons about Callings and the Church}
% -----------------------------------------------------------
	\label{LESSONS-CHURCH}
	
	% ----------------------
	\subsection{Working with YSAs}
	% ----------------------
	
		\begin{itemize}
		
			\item Lessons we learned from my time in the YSA bishopric
				\begin{itemize}
					\item We did not issue almost any callings after about the third week. We never returned to this in bishopric meeting, the result was that if you didn't get a calling right away, you never got one.
					\item Brett let the students fail. This was good. 
					\item We didn't give enough opportunities to the students to hang out in low-stakes smaller activities. 
					\item Brett wasn't very good at receiving feedback from us---there were times when we tried to be clear about what wasn't working and he either didn't listen or didn't care.
					\item Brett talked too much in meetings, like fifth Sunday meetings and sacrament meetings.
				\end{itemize}

			\item The most important thing in finding someone for a calling is finding who the Lord wants to be in that position.
				\begin{compactitem}
					\item Many times the Lord seems to be indifferent about this and there is no revelation to receive.
					\item But you must be open, even to very different ideas---remember Christian Clement, who was called to be EQ president in our YSA ward in fall 2022. He's weird, socially awkward, possibly or even likely gay, a bad public speaker, a poor leader, and a lot of people don't take him seriously in any way. But he did more toward fulfilling the calling and trying to minister to people in the ward in three weeks than both of our presidents in the two semesters prior to that.
				\end{compactitem}
		
		\end{itemize} 
	
	% ----------------------
	\subsection{Working with YSAs}
	% ----------------------
	
		\begin{itemize}
			
			\item \textbf{Lessons will be more effective if I understand that I should be trying to create a group therapy session, instead of trying to teach principles that I've been thinking about or studying.}
				\begin{compactitem}
					\item Principles are good, but what YSAs need is solidarity and support. They are going through many difficult things---I experienced those things too, but they are hard for me to remember.
					\item Don't plan a lesson with the intent of teaching a particular principle. I will lift them more if I try to create an environment where they can share about their difficulties and draw strength from each other.
					\item Along the way, I can teach principles and offer my own thoughts, but doing those things is not my primary aim.
				\end{compactitem}
				
			\item \textbf{Do not blind yourself by focusing on traditional ways of doing things, epsecially in the YSA ward. The goal is to help people, not do what has been done before, or what is traditional to do. Think of \textit{new ways} to do tihngs!}
			
			\item 
			
		\end{itemize}
		
% -----------------------------------------------------------
\section{Lessons about Fiction Writing} % *FICTION *WRITING
% -----------------------------------------------------------

	\begin{compactdesc}
	
		\item[\tr{001}] \textbf{In the standard Saves the Cat beat sheet structure, the parts of Act 2 before the midpoint are where you will do a lot of your characterizations of secondary characters, places, and events in order to get payoffs from these assets later on.}
			\begin{compactitem}
				\item This means that you have to fill the ``Fun and Games'' parts the secondary characters, so that readers learn about them through the POV characters.
				\item Think of Act 2 as moving pieces into place that you will then cash in on later as the things which have been moved into place begin to react and go off in the direction you want them to.
			\end{compactitem}
			
		\item[\tr{002}] \textbf{As you move forward in a book, you have to keep a character's motivations fresh and realistic, and sometimes you have to change them entirely in order to keep them moving where you ultimately want them to do.}
			\begin{compactitem}
				\item I learned this when I got to the midway point of \textit{Sword Without}. I realized that my original plan was to just give Rumara the very same motivations she had had previously, and I expected that to carry her all the way to Lakespring and to the end of the book.
				\item But I realized this was much too weak; she needed an entirely new motivation, and so I had to redo a few things in order to give her a new task to get her where she needed to go.
				\item The key here is discovering new subplots and stories that will keep the motivations new and interesting and powerful enough both for the characters and for the readers.
				\item You can't keep using the same motivations over and over again.
			\end{compactitem}
			
		\item[\tr{003}] \textbf{If you have outlined too much empty space in a story, and can't figure out how to get your character to where they need to be for their final destination, then consider bringing the final destination closer to them.}
		
		\item[\tr{004}] \textbf{After creating an initial outline, challenge yourself to shorten it by asking, what events can I cut? What events and chapters can I combine into one?}
		
		\item[\tr{005}] \textbf{You have to make something real and tangible happen in the first couple of pages. You get a grace period with the reader, but don't delay. Introduce the first few pieces, and make something real happen---the LAUNCH.}
			\begin{compactitem}
				\item The LAUNCH is the first real plot-driving event that happens.
				\item When the LAUNCH needs to take place depends on several factors, including genre and the amount of information the characters and reader have about the situation they are in and what is happening.
				\item For a book like \textit{Sword Without}, because the characters are fully aware of the situation they are in, you can have them speak openly about the LAUNCH, and the reader thereby comes to know about the LAUNCH fairly early.
					\begin{compactitem}
						\item This means that you have a bit more of a grace period with the setup, because the reader ALREADY KNOWS the LAUNCH is coming, and so has that to look forward to.
					\end{compactitem}
				\item For a book like \textit{Panopticon}, you've got Hans, who has no idea where he is, who the aliens are, or what is going on. 
					\begin{compactitem}
						\item This means that the LAUNCH needs to happen SOONER, since the reader has nothing else to go on or look forward to.
						\item The reader does not and cannot know about the LAUNCH, so you have to make plot-revelant stuff happen SOONER. The LAUNCH has to happen EARLIER.
						\item And I wonder if, in a book like this, the LAUNCH needs to be more dramatic, or more sudden, or something like that; it's going to come out of nowhere, and you want it to really grab.
						\item Consider also the utility of short chapters for a book like this.
					\end{compactitem}
			\end{compactitem}
			
		\item[\tr{006}] \textbf{In connection with 005, a major factor that must play a big role in structuring your plot is HOW MUCH THE CHARACTERS KNOW ABOUT THE SITUATION THEY ARE IN, because that level of knowledge determines HOW MUCH YOU WILL BE ABLE TO SHARE WITH THE READER ABOUT THE EVENTS THAT THEY OUGHT TO BE LOOKING FORWARD TO.}
		
		\item[\tr{007}] \textbf{If you are going to have an alien race or some other unfamiliar creature/animal play an important role in a story, you must give a DETAILED, CRYSTAL CLEAR description of what they look like.}
		
		\item[\tr{008}] \textbf{Things that need to happen in the setup:}
			\begin{compactitem}
				\item Hinting at the launch.
				\item If appropriate, hinting at the events surrounding the break into 2.
				\item Going with both of the previous things, giving the relationship between the launch and the break into 2.
			\end{compactitem}
			
		\item[\tr{009}] \textbf{The way I'm most creative and prolific in generating ideas for outlines:}
			\begin{compactitem}
				\item I learned this in working on the outline for \textit{Heaven's Key}.
				\item I floundered for weeks, maybe longer, trying to think of ideas for the third act and really feeling like I was getting nowhere.
				\item Then I had a few days in January and February 2024 when I was just sitting down to the mindmap and I started filling in details of characters and places.
				\item And then it just felt like new ideas were falling out of me, where I hadn't been able to have any new ideas before.
				\item \textbf{The lesson is: my best process involves just sitting down and building out the world and characters as much as possible---as I do this, my mind makes connections and comes up with tons of new ideas to use for the plot.}
				\item I can't force plot; I suck at just pulling plot ideas from nowhere.
				\item I'm much more creative when I'm building the world!
			\end{compactitem}
			
		\item[\tr{010}] \textbf{The only way to resolve mental blocks about the plot or other problems is to sit down and come up with other details for the story, including characters and their motivations.}
			\begin{compactitem}
				\item This is closely related to point \#9. 
				\item If you feel like the plot is a mess, it probably is.
				\item The way to fix it is to continue to develop other aspects of the story, including the setting, interesting event ideas, and so on.
				\item But the most important of all of these is the characters and their motivations.
				\item Once you start developing characters and giving them real, meaningful motivations, it will start to get easier to unravel plot problems.
			\end{compactitem}
			
		\item[\tr{011}] \textbf{The Circle Cycle of plot development}
			\begin{compactitem}
				\item I've been thinking about how getting ideas to work for \textit{Heaven's Key} has been so difficult. My struggle has led to the insights in numbers \#9 and \#10.
				\item But I think I've also arrived at a more general method: think of the main nodes in a mindmap like points on the outside of a circle. The nodes are things like places, characters, cool events, mysteries, plot points, questions, and so on.
				\item So the development process is basically going around in a circle, from main node to main node, continually deepening the points in each node, and as you do it, making connections and adding to other nodes.
				\item So you go to characters, and that gives you ideas for plots, and plots give you requirements for characters and places, and so on.
				\item You have the nodes in a circle, and you develop the story by cycling around the circle and adding material as you go.
			\end{compactitem}
			
		\item[\tr{012}] \textbf{Are you \textit{really} ready to start writing the book?}
			\begin{compactitem}
				\item In order to actually be ready, you need the following, in addition to all the story, setting, and character notes:
					\begin{compactitem}
						\item Maps of all key locations
						\item Lists of important characters or titles
						\item Lists of important vocabulary to remember to use (you add to this as you go along)
						\item Any other lists you need to remind of important things
						\item A long list of incidental characters that can be brought up or referred to
					\end{compactitem}
			\end{compactitem}
	
	\end{compactdesc}
	
% -----------------------------------------------------------
\pagebreak{}
% -----------------------------------------------------------